\documentclass{beamer}

\usepackage{amsmath}
\usepackage{amssymb}

\newcommand{\pro}{\,\textperiodcentered\,} % https://tex.stackexchange.com/a/631450/227665

\title{Diagonalization}
\subtitle{Pronounced: \textit{Di\pro{}a\pro{}go\pro{}na\pro{}li\pro{}za\pro{}tion}}
\author{Ahmed Mohammed}
\institute{Georgia Institute of Technology}
\date{PLUS Session, Summer 2023}

\usetheme{Berlin}
\usefonttheme{serif}
\usecolortheme{crane}

\begin{document}
	\begin{frame}
		\maketitle
	\end{frame}

	\section{Open Activity}
	\begin{frame}
		\frametitle{Open Activity: Motivation}
		\textbf{Question:} \pause
		\begin{center}
			If you know that a matrix has \(n\) eigenvectors that form a linearly independent set,\pause \ and I give you the corresponding eigen-stuff,\pause
			\ can you find the matrix?
		\end{center}
	\end{frame}

	\begin{frame}
		\frametitle{Open Activity: Introduction}
		\begin{definition}
			A matrix \(A \in \mathbb{R}^{n \times n}\) is diagonalizable if and only if it is \textbf{similar} to a diagonal matrix.
		\end{definition}
		
		By similar we mean that \(A\) could be expressed as \(A = PBP^{-1}\), where \(P\) is an invertible matrix and \(B\) is the matrix that \(A\)'s similar to.
		\pause

		If \(A\) is similar to a \emph{diagonal} matrix \(D\), then the matrix is diagonalizable. (Note that the entries of \(D\) are \emph{necessarily} the
		eigenvalues of \(A\)).
	\end{frame}

	\section{Main Activity: Exercises}
		\begin{frame}
			\frametitle{Main Activity: Why? (1)}

			Consider the earlier question.\pause\  Find the matrix whose eigenvalues are \(\lambda_1 = \lambda_2 = 2\) and \(\lambda_3 = 1\) with 
			corresponding eigenvectors \(\vec{v}_1 = \begin{bmatrix}1 \\ 0 \\ 1\end{bmatrix}\), \(\vec{v}_2 = \begin{bmatrix}2 \\ 1 \\ 0\end{bmatrix}\), 
			and \(\vec{v}_3 = \begin{bmatrix}1 \\ 1 \\ 0\end{bmatrix}\).
		\end{frame}
		
		\begin{frame}
			\frametitle{Main Activity: Why (2)}
			\begin{itemize}
				\item<1-> If \(A\) has \(n\) linearly independent eigenvectors, then so does \(A^T\).
				\item<2-> Is it possible that \(A\) is similar to a diagonalizable matrix \(B\), but \(A\) isn't diagonalizable?
				\item[\textbf{T/F}]<3-> Row-equivalent matrices are similar.
				\item[\textbf{T/F}]<4-> Every invertible matrix is diagonalizable.
				\item<5-> In the spirit of the previous question, find a \(3 \times 3\) matrix that is
						\begin{enumerate}
							\item<6-> Invertible and diagonalizable.
							\item<7-> Singular and diagonalizable.
							\item<8-> Invertible and not diagonalizable.
							\item<9-> Singular and not diagonalizable.
						\end{enumerate} 
			\end{itemize}
		\end{frame}
	\section{Closing Activity}
		\begin{frame}
			\frametitle{A Rabbit Hole}
			\begin{itemize}
				\item<1-> Diagonalization leads to very interesting ideas. If you want to know more, read Section 5.4, and go to office hours!
				\item<2-> Read the online Interative Linear Algebra textbook (ILA) (for more intuition)!
				\item<3-> So far, we've dealt with \emph{real} eigenvalues. What about if a real matrix has \emph{complex} eigenvalues?
			\end{itemize} \pause
			\onslide<4>{\[A = PCP^{-1}\] As it'll turn out, every real matrix that could be written in this form defines a rotation-dilation transformation!}
		\end{frame}
\end{document}
